% ============================================
% CONT-02
% Gerado de molambudos.md — build_miolo.py
% ============================================
\section*{\textbf{CONT-02} --- Você Estava Lá}
\markboth{CONT-02 --- Você Estava Lá}{}

\noindent \textit{Você sempre esteve}

\noindent Você acha que está lendo sobre o passado.

\noindent Acha que Joaquim é um personagem. Que a vala é uma cena. Que o olho amarelo é uma metáfora para o trauma. Que a criatura é uma projeção da mente dele, uma defesa psicológica, um mecanismo de dissociação.

\noindent Você está errado.

\noindent Você estava naquela vala.

\noindent Não fisicamente. Você não nasceu em 1907, não andou no sertão de Quixeramobim, não viu seu pai morrer de fome. Tudo isso é verdade.

\noindent Mas você estava lá de outro jeito.

\noindent Você estava lá porque a história de Joaquim é a história de todo mundo que já sentiu fome. Não fome de comida. Fome de algo que o mundo não dá. Fome de sentido. Fome de afeto. Fome de justiça. Fome de ser visto.

\noindent Você já sentiu essa fome.

\noindent Você já se sentiu sozinho numa multidão. Já teve a sensação de que há algo errado dentro de você. Algo que não tem nome. Algo que você alimenta com distrações, com trabalho, com barulho, com qualquer coisa que impeça o silêncio. Já acordou no meio da noite com a sensação de que não está sozinho no quarto.

\noindent Sinta o frio subir da nuca até o couro cabeludo. Seu estômago se aperta um segundo antes de você lembrar por quê.

\noindent A criatura que Joaquim encontrou na vala não é uma entidade sobrenatural.

\noindent Ela é a fome que todo mundo carrega. A diferença é que Joaquim \textit{viu} ela. Tocou nela. Deixou ela entrar.

\noindent Ela já está em você também.

\noindent Você só não viu ainda.

\noindent Mas depois deste livro, você vai começar a ver.

\noindent Nos cantos do quarto.

\noindent No reflexo do vidro.

\noindent No olho amarelo do cachorro na rua.

\noindent E numa noite, quando você estiver deitado, com a luz apagada, você vai sentir um peso na nuca. Um hálito quente.

\noindent Não se mexa.

\noindent Não vai ter medo. Você vai reconhecer.

\noindent Você vai pensar: \textit{\textquotedblleft{}Você também esteve lá.\textquotedblright{}}

\noindent E a criatura vai responder:

\noindent \textit{\textquotedblleft{}Nunca saí.\textquotedblright{}}

\vspace{0.5em}
\noindent\textit{\textbf{↪ Rotas:} \rota{CONT-03} \rotasep \rota{MEM-04}}